\section{电流与电流密度}
电荷的运动将在电场之外产生磁场，而大量定向运动的电荷将形成电流，因此在开始本章对磁场的研究前，我们有必要先对电流以及如何定量描述电流分布做一些更多的探讨。

电荷在电场作用下的定向运动形成\uwave{电流}，而形成电流的带电粒子统称为\uwave{载流子}，最为常见的载流子是导体中的电子。但除此之外，在溶液中的载流子还可以是\uwave{阴阳离子}，在半导体中的载流子还可以是带正电的\uwave{空穴}。接下来电流的讨论均以金属导体中的电子为背景展开。

电场作用下，正负电荷所受电场力相反，定向运动的方向也相反，其中我们将正电荷运动的方向规定为电流的方向。例如，虽然在金属导体中的载流子实际是带负电荷的电子，但是我们就可以将等效视为沿反向运动的等量正电荷，此时电流的方向就是电子运动的反方向。\footnote{由此可见，电流并不是载流子的运动方向，而是正载流子的方向或负载流子的反方向。}

由此可见，产生电流的条件是
\begin{enumerate}
    \item 存在可以自由移动的电荷。
    \item 存在电场的作用。
\end{enumerate}
这也就解释了为何电介质是绝缘体，电介质中没有可以自由移动的电荷，也就无法导电了。

现在的问题是，如何定量的描述电流？

\subsection{电流强度}
\uwave{电流强度}是一个我们所熟知的一个用于描绘电流的物理量\footnote{在不引起误会时，电流强度也可以直接简称为电流。}，它反应了电流的强弱，其被定义为单位时间内通过导体某一横截面的电荷量，通常用符号$I$表示，即
\begin{BoxDefinition}[电流强度]
    I=\Fdif{Q}{t}
\end{BoxDefinition}
电流强度在国际单位制中的单位是\si{A=C\cdot s^{-1}}，称为\uwave{安培}，量纲是\si{[I]}。

电流强度是标量，通常所说的电流强度的方向是指电流的流向，即电流是经由界面自左向右流动还是自右向左流动，有时我们也用电流强度的正负来标记电流的流向。

如果我们研究和考虑的情形比较简单，例如电流在一根导线中流动，这根导线的粗细和材质，或许是均匀的，或许在局部有些不均匀，但是如果我们只需要知道单位时间内留过这导线某一截面的电荷量，并且并不关心在那截面上细微而具体的电流分布，那这些（粗细和材质的均匀与否）都不是很重要，电流强度已足矣反映我们关于电流需要知道的一切了。

但是若我们想考虑的更细致一些，如\Refx{图}{电流在粗细不均匀的导线中的流动}所示，很明显，导线较粗处电流也较为稀疏，导线较细处电流就较为密集，但是通过任何一处截面的电流强度是完全相同，那么如果我们想定量的研究一些这种“电流的疏密”，此时电流强度就无能为力了。这尚且还是在导线中的情况，这里电流强度至少还能描述一些东西。

可是在实际中我们还会遇到电流在大块异形导体中流动的情况，例如通过向大地通电流进行地质勘探。此时导体中电流的方向处处都不同，其根本没有一个基本一致的流向，其是完全散开的，这时电流强度就完全无法反映任何信息了，我们甚至都不知道该取哪一截面！

\newpage

\begin{Figure}{电流在粗细不均匀的导线中的流动}
    \inputfigure[scale=0.6]{Chapter10A_Figure01}
\end{Figure}

\subsection{电流密度}
由以上讨论可知，我们有必要在电流强度外定义一个描述电流分布的物理量，这应该是一个空间矢量函数，其方向反映该处电流的流向，其大小反映该处的电流的强弱，很明显，若导体中某一点处的电荷越多且电荷运动速度越快，那么该点处的电流就越强。

由于导体中的电荷是体分布的，因此我们就将导体中某一点处的电荷体密度$\rho$与该点处的电荷运动速度$\vb*{v}$的乘积定义为该点处的\uwave{电流密度}，通常用符号$\vb*{J}$表示，即
\begin{BoxDefinition}[电流密度]
    \vb*{J}=\rho\vb*{v}
\end{BoxDefinition}
电流密度在国际单位制中的单位是\si{A\cdot m^{-2}=C\cdot m^{-2}\cdot s^{-1}}，量纲是\si{[L^{-2}I]}。

通过\Refx{定义}{电流密度}，我们就可以更好理解什么是电流强度了，电流密度描述了电流的分布，电流强度反映的是通过一个截面的总电流量，即\textbf{电流强度是电流密度的通量}。
\begin{BoxEquation}[电流与电流密度]
    I=\Isnt[S]{\vb*{J}\cdot\dif{\vb*{S}}}
\end{BoxEquation}
我们可以用\Refx{图}{电流在粗细不均匀的导线中的流动}中那样的流线图将电流场形象的表现出来，称为\uwave{电流线图}，其疏密反映电流密度的大小，其方向反映电流密度的方向，如果我们在电流线图中取一个截面，那么通过该截面的电流线数目即该截面上的电流强度。

我们现在来看如果所取的“截面”是一个闭合曲面时会发生什么，电流密度是等效正电荷的流动方向，电流此时表示的应是单位时间内流出闭合曲面的电荷量，而根据\Refx{定律}{电荷守恒定律}，电荷不可能凭空出现或消失，因此闭合曲面内的的总电荷量$q$的变化完全来自电荷流入或流出该闭合曲面，故此时\Refx{公式}{电流与电流密度}就可以表示为
\begin{BoxPrinciple}[电流连续性方程]
    流出闭合曲面的电流，即该闭合曲面内的电荷量的负变化率。
    \begin{BoxEq}
        I=\Isot[S]{\vb*{J}\cdot\dif{\vb*{S}}}=-\Fdif{q}{t}
    \end{BoxEq}
\end{BoxPrinciple}
这里的负号是因为电流强度的正值表示电荷向闭合曲面外运动，即为电荷的负变化率，这与先前的\Refx{定义}{电流强度}并不矛盾，那时的电荷量是指流经截面的电荷，意义不同。

这里我们也可以将\Refx{定律}{电流连续性方程}以微分形式表示
\begin{BoxPrinciple}[电流连续性方程的微分形式]
    电流密度的散度，即该处电荷密度的负变化率。
    \begin{BoxEq}
        \nabla\cdot\vb*{J}=-\Fdif{\rho}{t}
    \end{BoxEq}
\end{BoxPrinciple}
这就说明电流的源完全来自电荷的增加与减少
\begin{itemize}
    \item 如果我们在某处看到一个电流的正源，那么该处的电荷一定在不断减少。
    \item 如果我们在某处看到一个电流的负源，那么该处的电荷一定在不断增加。
\end{itemize}\vspace{-15pt}

\subsection{恒定电流场与恒定电场}
现在让我们来考虑一种特殊情况，假设电流密度$\vb*{J}$是与时间无关的，此时电流分布是恒定不变的，称为\uwave{恒定电流}，而我们知道，电流是电荷在电场作用下的定向运动，电流分布不变，电场分布就不变，电荷分布也就不变，即$\rho$和任意闭曲面内的$q$是与时间无关的常数，故
\begin{BoxPrinciple}[电流恒定条件]
    恒定电流场在闭曲面上的通量为零，流入闭曲面的电流总是等于流出闭曲面的电流。
    \begin{BoxEq}
        I=\Isot[S]{\vb*{J}\cdot\dif{\vb*{S}}}=0
    \end{BoxEq}
\end{BoxPrinciple}
\begin{BoxPrinciple}[电流恒定条件的微分形式]
    恒定电流场的散度处处为零，电流分布与电荷分布不随时间变化。
    \begin{BoxEq}
        \nabla\cdot\vb*{J}=0
    \end{BoxEq}
\end{BoxPrinciple}
这就是说，恒定电流场必定是无散的，虽然电荷在导体中在不断流动，但是电荷的分布却没有发生任何变化，从而产生的电场也是恒定不随时间变化的，称之为\uwave{恒定电场}。

恒定电场与静电场具有相似之处，两者的电荷分布和电场分布都不随时间变化，因此描述静电场基本性质的\Refx{定律}{电场高斯定理}和\Refx{定律}{电场环路定理}对于恒定电场也是完全适用的，因此我们也可以在在恒定电场中也可以引入\Refx{定义}{电势}的概念。

恒定电场与静电场最主要的区别在于达成电荷分布不随时间变化的方式，后者的电荷是真正静止的，前者则是通过电荷在定向运动中的动态平衡，即某点处在有电荷流出时必然有等量电荷流入进行补充，从而实现电荷分布不变。因此，静电场中的静电平衡理论，导体内的场强处处为零，导体表面等势，在恒定电场中就不再适用了，在恒定电场中的导体
\begin{itemize}
    \item 导体表面具有一定的电势差，且不随时间变化。
    \item 导体内部具有恒定的电场强度，因而能使电荷定向运动形成恒定电流。
\end{itemize}
恒定电场最直观的一个例子就是电路，导线虽然是一种导体，但在接入电源后，导线却并不是等势的，导线正负两端的电势是不同的。这也反映出了恒定电场的另外一个特征，即恒定电场的维持必定是需要能量的，例如电源。相比之下，静电场的维持不需要能量。

\subsection{恒定电路与欧姆定律}
恒定电流场需要电荷源源不断的流动，这种情形只有在那种首尾相接的路线中才有有可能出现，这种由导线构成的完整回路称为\uwave{电路}，当然，电路并不总是能形成恒定电流场，电路中可以形成恒定电流场的称为\uwave{恒定电路}，电阻电路是恒定电路，电容电路则不是。

% 恒定电路最大的特点是其电学参量是不随时间变化的常量，对于电阻电路，在开关闭合后，电压和电流便是我们计算的那个数字了，现在如此，将来也如此。相反，对于电容电路，正如先前我们推导的\Refx{公式}{电容充电过程的电压}和\Refx{公式}{电容放电过程的电压}所示的那样，电压和电流都是随时间变化的函数，这种非恒定电路的情况就较为复杂了。

恒定电路的一个重要特点是，\textbf{在同一支路上的电流处处相等}，证明是非常容易的，我们可以在导线上取两个截面，连同其间的导线表面取作一个完整的闭合曲面，而后者上显然不会通过任何电流，结合\Refx{定律}{电流恒定条件}可知，从后截面流入的电流总是等于从前截面流出的电流，由于这样的截面是可以任选的，因此任意截面上的电流都是相同的。

% 这一结论是非常重要的，虽然我们使用的导线通常是均匀的，但是导线间连接的电阻器的截面积与导线的截面积必然是不同的，只有通过这一结论，我们才可以坚定的说
% \begin{center}
%     “瞧！电阻上的电流与导线相同！”
% \end{center}
恒定电路中的电流是常数，而实验表明，在大部分情况下，电路两端的电压总是与电流成正比，这一比例系数称为\uwave{电阻}，记作$R$，这一经验结论就是我们非常熟悉的\uwave{欧姆定律}
\begin{BoxPrinciple}[欧姆定律]
    电路两端的电压总是与该电路上的电流成正比。
    \begin{BoxEq}
        U=IR
    \end{BoxEq}
\end{BoxPrinciple}
电阻在国际单位制中的单位是~\si{\ohm=V\cdot A^{-1}}，称为\uwave{欧姆}，量纲是\si{[L^{2}MT^{-3}I^{-2}]}。

电阻的倒数称为\uwave{电导}，记作$G$\vspace{-5pt}
\begin{BoxDefinition}[电导]
    G=\frac{1}{R}
\end{BoxDefinition}
电导在国际单位制中的单位是\si{S=\ohm^{-1}=A\cdot V^{-1}}，称为\uwave{西门子}，量纲是\si{[L^{-2}M^{-1}T^3I^2]}。

应当指出，\Refx{定律}{欧姆定律}并不总是成立的，使得欧姆定律成立的器件称为\uwave{欧姆元件}，否则称为\uwave{非欧姆元件}，其中一个我们熟悉的例子是小灯泡，它的电阻会随着温度变化。\footnote{小灯泡本质仍然是电阻丝，只不过为了发光，在电阻丝上的温度变化范围远超正常的电阻器，以至于电阻率有了明显变化。}

导体电阻的大小与导体的材料和几何形状有关，对于由一定材料制成的横截面均匀的导体，电阻与导体长度$L$成正比，电阻与导体横截面积$S$成反比，这一结论即\uwave{电阻定律}
\begin{BoxEquation}[电阻定律]
    R=\rho\frac{L}{S}
\end{BoxEquation}
上式中的$\rho$不是电荷密度，其表示材料的\uwave{电阻率}，其倒数称为\uwave{电导率}。
\begin{BoxDefinition}[电导率]
    \gamma=\frac{1}{\rho}
\end{BoxDefinition}
电阻率在国际单位制中的单位是\si{\ohm\cdot m^{-1}}，量纲是\si{[LMT^{-3}I^{-2}]}。

电导率在国际单位制中的单位是\si{S\cdot m}，量纲是\si{[L^{-1}M^{-1}T^{-3}I^{-2}]}。

现在让我们来研究一下\Refx{定律}{欧姆定律}在微观上会导致什么样的结果，我们在导体中取一个很小的\uwave{电流管}，即沿电场线的一个管状假象曲面，由于这里取的电流管很小，可以近似视为一个柱体，设电流管的长度为$\dif{L}$，设电流管的截面积为$\dif{S}$，那么
\begin{Equation}[欧姆定律的微分形式1]
    \dif{I}\xlongequal{\text{\Refx{定律}{欧姆定律}}}\frac{\dif{U}}{R}
    \xlongequal[\text{\Refx{定义}{电导率}}]{\text{\Refx{公式}{电阻定律}}}
    \gamma\dif{U}\frac{\dif{S}}{\dif{L}}
\end{Equation}
其中$\dif{I},\dif{U}$分别是小电流管上的电流和其两端的电压。

由于小电流管中的电场可以近似视为匀强电场，因而
\begin{Equation}[欧姆定律的微分形式2]
    \dif{U}=E\dif{L}
\end{Equation}
将式\ref{式:欧姆定律的微分形式2}代入\ref{式:欧姆定律的微分形式1}，消去$\dif{L}$
\begin{Equation}[欧姆定律的微分形式3]
    \dif{I}=\gamma E\dif{S}
\end{Equation}
由于小电流管截面上的电流可以近似认为是均匀分布，故有$\dif{I}=J\dif{S}$，代入上式约去$\dif{S}$
\begin{equation*}
    J=\gamma E
\end{equation*}
考虑到电流密度的方向应与场强保持一致
\begin{BoxGeneral}[欧姆定律的微分形式][定律]
    \begin{BoxEq}
        \vb*{J}=\gamma\vb*{E}
    \end{BoxEq}
\end{BoxGeneral}
值得注意的是，上式虽然是由恒定电路的情形导出，但是对于那些非恒定电路，由于\Refx{定律}{欧姆定律}在电阻两端仍然是成立的，上式实际上也是适用的\footnote{只不过此时$\vb*{J},\vb*{E}$都会随时间变化，但是两者间的正比关系$\vb*{J}=\gamma\vb*{E}$始终成立，就像$U=IR$在$U,I$变化时也成立。}。该结论在电磁场中的地位与上一章中\Refx{公式}{电位移矢量与场强的电容率表示}是相当的，但同样，这也是对物质特性的一种近似，对于那些非欧姆元件，这一结论就未必正确了。我们仍然以小灯泡为例来说明，如果认为决定电阻的温度只与电流有关，那么此时$\vb*{J},\vb*{E}$是非线性的，如果我们还要考虑温度在小灯泡上的积蓄，那么此时$\vb*{J},\vb*{E}$之间就不存在简单的单值对应关系了。